\scp{A mixed metaphor}{02-07}
The issues above reflect only a sample of things that have added
confusion rather than clarity in unravelling the historical-comparative
picture of the region.
Most of the arguments about linguistic prehistory
and subgrouping in the region to date have been based on very little data.

In looking at Austronesian speaking peoples speaking mostly typologically
SVO languages coming into a relatively large
and diverse region already populated by people speaking various
typologically SOV Papuan languages,
we have a picture something like a river of fresh water coming
into contact with a body of salt water (see \frf{02-28}).
There are places where the salt water is demarcated
and is clearly salt water,
and the fresh water is clearly fresh water.
And places where it is something different altogether -- it is brackish.
The fresh water is not uniform where it comes into the sea,
with a deep channel on the side where the water is running swiftly,
shallower and running more slowly on the opposite side,
with multiple tributaries at the mouth flowing around small islands,
large rocks and sandbars.
The sea is also not uniform with different temperatures
and ecosystems at different distances from shore
and at different depths.

\begin{figure}
	\centering
	\includegraphics[width=0.45\columnwidth]{02-28.jpg}
	\fcp{Wae Tina Nangan river mouth in south Buru}{02-28}
\end{figure}

Once the waters meet,
there are places where the identity as one or the other is not
always clear, because they have mixed; they are no longer purely fresh water
or purely sea water.
They start looking a bit the same in places.
But the mixture is not the same throughout the zone of contact.
Different bottom features (substrate),
channels, islands, eddys and back currents mean that there is
great variety within the area where the two meet.
And silt and debris floating around or stirred up by the two
currents coming together add distractions
and dangers to the unwary.
And there is not just one river flowing into the sea.
For example, not far to the \ili{east} of the river mouth shown in
\frf{02-28} is another small river,
Namrole (Role Bay -- that gives the large bay its name),
flowing into the same sea.
And not far to the west of the picture is Wamkana (Forceful River),
a small river flowing with great force.

Experienced boatmen know they must be extra careful where rivers
flow into the sea,
because of the unpatterned waves
and currents, and unexpected debris that can flip or damage a boat.
They also know that the intersection of the two bodies of flowing
water means the hidden bottom is ever-changing
and a bit unpredictable -- not the same as the last time they passed by.
